\documentclass[a4paper,11pt]{article}
% \usepackage[latin1]{inputenc}
\usepackage[utf8]{inputenc}
\usepackage[brazil]{babel}
\usepackage{float,graphicx,graphics,amssymb,amsfonts,newlfont,indentfirst}
\usepackage[centertags]{amsmath}
\usepackage{fancyhdr}
\usepackage{ragged2e}
\usepackage{geometry}
\geometry{top=0cm,bottom=1cm,left=3cm,right=2cm}
\usepackage[normalem]{ulem}
\pagestyle{fancy}
\usepackage{lineno} %pacote para colocar números das linhas(será utilizado pelos pareceristas, tirar na versão final%
\usepackage{url}
\usepackage{color}
%\usepackage{quoting}
 \fancyfoot{} %tira número das páginas do rodapé
%%%%%%%%%%% Fonte padrão - ptm (Times) %%%%%%%%%%% USAR ESTA FONTE
\renewcommand{\rmdefault}{ptm}
\renewcommand{\sfdefault}{ptm}
\renewcommand{\ttdefault}{ptm}
%%%%%%%%%%%%%%%%%%%%%%%%%%%%%%%%%%%%%%%%%%%%%%%%%%

\renewcommand{\headrulewidth}{0pt} % NÃO ALTERAR, retira linha automática de cima

\usepackage{mathptmx} % Usar esta Fonte padrão em fórmulas - Times!!!!

\headheight 30mm      %
\oddsidemargin 2.0mm  %
\evensidemargin 2.0mm %
\topmargin -18mm       %
\textheight 235mm     %
\textwidth 160mm      %
\headsep 2.5mm        %
\parindent 1mm        %

%\lhead{\epsfig{file=fig0.eps, height=2.5cm, width=16cm}}
\chead{\includegraphics[height=2cm, width=16.2cm]{bannerroxo.png} }


\begin{document}
\linenumbers %coloca números nas linhas a partir daqui (TIRAR NA VERSÃO FINAL)
%\vspace{20cm}
\centering{{\Large{\bf  \vspace{5cm}Instruções para preparação de trabalhos a serem apresentados no ERMAC 2024}}\\{\small \color{red} (O título deve ser em negrito, com espaçamento simples e centralizado. A primeira letra da primeira palavra deve ser maiúscula, e as demais, minúsculas, exceto para nomes próprios. O subtítulo, se houver, deve ser com letra minúscula e sem negrito. Observações como esta, em vermelho, devem ser excluídas ao elaborar o resumo.)}}

\begin{flushright}{\it
\vspace*{5mm}
Fulano Escritor {\color{red}(nome completo do 1º autor, conforme usado em referências)}\\
Universidade; Faculdade\\
xxxx@yyy.zzz {\color{red}(e-mail do 1º autor, se possível, institucional)}

\vspace*{5mm}
\underline{Beltrano Falante }{\color{red}(nome completo do 2º autor, conforme usado em referências)}\\ %neste caso o segundo autor é o apresentador
Universidade; Faculdade\\
aaaa@bbbbb.ccc {\color{red}(e-mail do 2º autor, se possível, institucional)}

\vspace*{5mm}
Sicrana Pensadora {\color{red}(nome completo do 3º autor, conforme usado em referências)}\\
Universidade; Faculdade\\
nnnn@ppppp.qqq {\color{red}(e-mail do 3º autor, se possível, institucional)}
}\end{flushright}


\setcounter{equation}{0}



\justifying{
{\noindent \bf Resumo:}{\small{ {\color{red}\it(desconsiderar este item para a Categoria 1)} O propósito deste template é apresentar o formato a ser utilizado para os trabalhos a serem submetidos para apreciação da Comissão Científica do ERMAC 2024, Bauru, SP. Estão previstas duas categorias, 1 e 2, onde os autores poderão submeter seus trabalhos. {\bf O resumo, necessário apenas para a Categoria 2}, deve conter, no máximo, 200 palavras. Ele deve ser escrito em fonte Times New Roman (já no preâmbulo), small, e deve enunciar de forma clara e sintética: problema de pesquisa, abordagem metodológica empreendida, resultados e conclusões. \\


}}
\vskip 0.2cm

{\small{{
\noindent{\bf{Palavras-chave:}} Matemática Discreta. Tableaux. Lógica RM3.\\
{\color{red} (Usar, no máximo, 4 palavras-chave, em fonte Times New Roman, tamanho 10, espaçamento simples, justificado. A letra inicial de cada palavra deve ser escrita em letra maiúscula, exceto as preposições, advérbios, conjunções, etc, e separadas por ponto final. Pelo menos uma delas deve corresponder a um dos eixos temáticos do evento: 01.Análise Aplicada, 02.Biomatemática, 03.Computação Científica, 04.Computação Gráfica, 05.Controle e Teoria de Sistemas, 06.Matemática Aplicada à Economia e a Finanças, 07.Matemática Aplicada à Engenharia, 08.Matemática Aplicada à Física, 09.Matemática Discreta, 10.Mecânica dos Fluidos e Aplicações, 11.Métodos Estocásticos e Estatísticos, 12.Métodos Numéricos e Aplicações, 13.Modelagem Matemática e Aplicações, 14.Otimi\-zação, 15.Problemas Inversos, 16.Processamento de Sinais, 17.Ensino, 18.Cálculo Fracionário e Aplicações.)
}}}}}

\section*{Introdução}
\quad
\justifying
 Os trabalhos aprovados para apresentação no ERMAC 2024 devem ser escritos, preferencialmente, na língua portuguesa, eventualmente, em língua inglesa, seguindo este template. Atenção: os trabalhos submetidos deverão estar rigorosamente num dos formatos disponíveis no site (MS Word ou LaTeX).


\section*{Submissão}
\quad
\justifying
Para a {\bf Categoria 1}, o trabalho deve ter, no máximo, {\bf duas páginas} e, se aceito, será apresentado por meio de um painel. Para a {\bf Categoria 2}, o trabalho deve ter de  {\bf cinco a sete} páginas e, se aceito, será apresentado oralmente. O nome do autor apresentador deve ser \underline {sublinhado}, para auxiliar a sua identificação pelos congressistas.

As Equações 1 e 2 são exemplos de formatação adequada para formas longas e curtas, respectivamente. A Equação de Navier-Stokes, juntamente com a equação de continuidade, em forma adimensional, são dadas por:

\small
\begin{eqnarray}
\nonumber
\frac{\partial {\mathbf u}}{\partial t} + \nabla \cdot ({\mathbf uu}) &=& \frac{1}{\rho}\left\{-\nabla p + \frac{1}{Re}\,[\nabla \cdot (2 \mu S)] +\right.
\\
&& \left.+\frac{1}{Fr^2}\,\rho {\mathbf g} +
\frac{1}{We}\,\kappa \delta ({\mathbf x} - {\mathbf x}_f) {\mathbf n}\right\},
%\text{,}
\label{eq1}
\end{eqnarray}
e
\begin{equation}
\nabla \cdot {\mathbf u} = 0,
%\text{,}
\label{eq2}
\end{equation}
\normalsize
sendo $Re = \rho_{0}UL/ \mu_{0}$, $Fr = U/ \sqrt{Lg}$ e $We = \rho L U^2 /\sigma_{0}$ os números de Reynolds, Froude e Weber, respectivamente.

Tabelas, figuras e equações devem ser referenciadas com a enumeração em algarismos arábicos. Por exemplo, a Equação (\ref{eq1}) apresenta uma expressão longa em duas linhas, a Tabela \ref{fonts} indica os formatos de texto das diferentes partes do documento e a Figura \ref{niveis} mostra que os gráficos podem ser coloridos.


\begin{table}[h]
\caption{{\small Tipos de tamanhos de letra nas partes deste documento.}}\label{fonts}
\centering
\begin{tabular}{|c|c|c|c|}
\hline
{\bf Texto} & {\bf MS Word} & {\bf LaTeX} & {\bf Aparência} \\
\hline
título & 14pt & Large & \textbf{bold}\\
autor(es), instituição, e-mail & 10pt & small & \textit{Itálico}\\
resumo, palavras-chave  &	10 pt	 & small  &	{\small normal}\\
texto principal	 &11 pt & normal & normal\\
\hline
\end{tabular}
\end{table}

O texto de legenda, para as tabelas e figuras, deve descrever os elementos principais das mesmas. Em LaTeX com figuras .jpg, usar pdflatex.

\begin{figure}[h]
\centering \includegraphics[height=4.5cm]{fig.jpg}\\
\caption{{\small Curvas de níveis da função $f(x,y)=x^3-3 x y^2$. }}
\label{niveis}
\end{figure}

{\bf O arquivo contendo o resumo finalizado deve ser submetido em formato PDF. O nome do arquivo deve conter três termos separados pelo símbolo underline ( \_ ), a saber: o número do eixo temático em dois dígitos, o primeiro nome do primeiro autor e o último sobrenome deste (por exemplo, 01\_Paulo\_Silva.pdf).}

\section*{Seleção de trabalhos}
\quad
\justifying
 Os trabalhos submetidos, dentro do prazo estabelecido, serão enviados aos revisores da Comissão Científica do ERMAC 2024. Com base nos pareceres da comissão, o trabalho poderá ser (1) aceito plenamente, (2) aceito sob a condição de que correções menores sejam feitas em curto prazo ou (3) rejeitado.


\section*{Citações}
Devem seguir as normas da ABNT NBR 10520. Nas citações, as chamadas pelo sobrenome do autor devem ser em letras maiúsculas e minúsculas e, quando estiverem entre parênteses, devem ser em letras maiúsculas.

Exemplos:

A ironia seria assim uma forma implícita de heterogeneidade mostrada, conforme a classificação proposta por Authier-Reiriz (1982).

\vskip 0.2cm
``Apesar das aparências, a desconstrução do logocentrismo não é uma psicanálise da filosofia [...]'' (DERRIDA, 1967, p. 293).

\vskip 0.2cm
a) As citações diretas, no texto, com mais de três linhas, devem ser destacadas com recuo de 4 cm da margem esquerda, espaço entre linhas simples e sem aspas, em fonte Times New Roman, tamanho 10.


\vskip 0.2cm
b) As citações diretas, no texto, de até três linhas, devem ser escritas entre ``aspas'' duplas e incorporadas ao texto.
Exemplos:
Barbour (1971, p. 35) descreve: ``O estudo da morfologia dos terrenos [...] ativos [...]''

\vskip 0.2cm
``Não se mova, faça de conta que está morta.'' (CLARAC; BONNIN, 1985, p. 72).

\vskip 0.2cm
Segundo Sá (1995, p. 27): ``[...] por meio da mesma `arte de conversação' que abrange tão extensa e significativa parte da nossa existência cotidiana [...]''

\vskip 0.2cm
c) Nas citações diretas, especificar no texto o ano de publicação e a(s) página(s) da fonte consultada. Estes dados devem ser colocados entre parênteses e separados por vírgula. Nas citações indiretas, a indicação da(s) página(s) consultada(s) é opcional, mas o ano de publicação da obra é obrigatório e deve estar entre parênteses.

\section*{Notas de rodapé}
Devem ser evitadas. Porém, caso necessárias, devem ser sintéticas e reduzidas ao máximo. Podem vir ao final da página, numeradas em sequência, em fonte Times New Roman, alinhamento justificado e espaçamento simples.

\vskip 0.2cm
\section*{Conclusões}
Aqui devem ser apresentadas as principais conclusões do trabalho, com base nas seções anteriores.

\section*{Submissão para a revista C.Q.D.}

Os trabalhos apresentados no ERMAC 2024 serão em breve publicados, em forma digital, no Caderno de trabalhos completos e resumos ERMAC 2024. Lembramos que os autores podem submeter trabalhos, até 31/8/24, para uma edição especial da C.Q.D.-- Revista Eletrônica Paulista de Matemática. 

As novas submissões devem basear-se nos trabalhos apresentados no evento, sendo possível e recomendável a realização de mudanças apreciáveis de conteúdo, com maior riqueza de resultados, solidez de argumentos e qualidade de apresentação.
Além disso, é necessário citar o trabalho publicado no Caderno de trabalhos completos e resumos e explicar brevemente, no texto principal, as mudanças realizadas.
As normas da revista estão disponíveis no link \url{https://sistemas.fc.unesp.br/ojs/index.php/revistacqd/about/submissions} 

\section*{Agradecimentos}
Os autores podem apresentar os agradecimentos a pessoas e instituições. Por exemplo, a Comissão Editorial do ERMAC 2024 agradece aos autores por seguirem as recomendações descritas neste template.

\section*{Referências}
A bibliografia deverá seguir o padrão da ABNT NBR 6023, separadas entre si por uma linha em branco, estar em ordem alfabética pelo sobrenome do primeiro autor, se necessário, usando-se, ainda, ordem cronológica, para trabalhos de um mesmo autor. Trabalhos dos mesmos autores, publicados no mesmo ano, devem ser listados utilizando-se a ordem alfabética do título do trabalho. Basicamente, as referências devem conter as iniciais dos nomes dos autores, sendo escrito, por extenso, apenas o último sobrenome. Seguem alguns exemplos:

\vskip 0.5cm
\begin{flushleft}

{\color{red} Livro com até 3 autores:}

GAUTSCHI, W. A survey of gauss-christoffel quadrature formulae. In: BUTZER, P.L.; FEHER, F. (Edit.). {\bf E. B. Christoffel: the influence of his work on mathematics and the physical sciences}. Basel; Boston: Birkhauser Verlag, 1981. p. 72-147.

\vskip 0.2cm
BRUNETTI, F. \textbf{Mecânica dos fluidos}. 2.ed. São Paulo:Pearson, 2005.\\
\vskip 0.2cm
JAIN, A.; ROSS, A.; NANDAKUMAR, K. \textbf{Introduction to Biometrics}. 1. ed. New York: Springer, 2001.

\vskip 0.4cm
{\color{red}Livro com 4 autores ou mais:}

ARENALES, M. et al. \textbf{Pesquisa Operacional}. Rio de Janeiro: Elsevier, 2006.

\vskip 0.4cm
{\color{red}Artigo em periódico:}

AVILA, A. Density of positive Lyapunov Exponents for $SL(2,\mathbb{R})$ cocycles. \textbf{Journal of the America Mathematical Society}, v. 24, n.4, p. 999-1014, 2011.

\vskip 0.4cm

KURODA, L. K. B. et al. Método da transformada diferencial generalizada no modelo fracionário de Malthus. \textbf{C.Q.D. -- Revista Eletrônica Paulista de Matemática}, Bauru, v. 10, p. 68-78, dez. 2017. Edição Ermac. 

\vskip 0.4cm
{\color{red}Trabalho em evento:}

BRAYNER, A. R. A.; MEDEIROS, C. B. Incorporação do tempo em SGBD orientado a objetos. In: SIMPÓSIO BRASILEIRO DE BANCO DE DADOS, 9., 1994, São Paulo. \textbf{Anais...} São Paulo: USP, 1994. p. 16-29.

\vskip 0.4cm
{\color{red}Dissertações e teses:}

CHERRI, A.C. \textbf{Reaprendendo tópicos de cálculo diferencial com o auxilio de softwares matemáticos}. 2001. 154 f. Trabalho de Conclusão de Curso (Especialização) \textendash\ Faculdade de Ciências, UNESP, Bauru, 2001.\\
\vskip 0.2cm
DINIZ, G. L. \textbf{A mudança no habitat de populações de peixes:} de rio a represa -- o modelo
matemático. 1994. 90 f. Dissertação (Mestrado em Matemática Aplicada) \textendash\ Unicamp, Campinas, 1994.\\
\vskip 0.2cm
OLIVEIRA, E. L. \textbf{Torres de Extensões Abelianas
de grau primo ímpar não ramificado}.2015. 63f. Tese (Doutorado em Matemática) \textendash\ IBILCE/UNESP, São José do Rio Preto, 2015.
\vskip 0.4cm
{\color{red}Homepages:}

INTERNATIONAL GEOGEBRA INSTITUTE. Matemática dinâmica para se aprender e se ensinar. 2014. Disponível em: http://www.geogebra.org/cms. Acesso em: 17 dez. 2014.
\end{flushleft}

\end{document}

